\section{Arquitetura do sistema}

O sistema foi organizado em quatro serviços principais sob Docker Compose.
Os três nós compartilham o mesmo código-fonte e mudam apenas por configuração de ambiente.

\begin{table}[h]
    \centering
    \begin{tabularx}{\textwidth}{@{} l l X @{} }
        \toprule
        Componente & Papel & Observações \\
        \midrule
        \code{node1}, \code{node2}, \code{node3} & Nós distribuídos & Executam FastAPI, Uvicorn e HTTPX; cada um mantém estado local em memória. \\
        \code{resource} & Observador & Registra entrada, saída e sobreposição na seção crítica; não decide prioridade. \\
        \code{docker-compose.yml} & Orquestração & Sobe a pilha local com um contêiner por nó. \\
        \code{Dockerfile} & Empacotamento & Define a imagem base da aplicação. \\
        \code{scripts/} & Automação & Contém smoke tests, demonstração e experimentos. \\
        \code{tests/} & Validação & Suíte unitária para relógio, mutex, eleição e estado. \\
        \bottomrule
    \end{tabularx}
    \caption{Componentes principais do projeto}
\end{table}

A comunicação entre os processos usa envelopes JSON enviados por HTTP.
Essa decisão preserva a visibilidade do fluxo de mensagens e facilita a inspeção por logs e por scripts de validação.
Cada envio aparece como chamada HTTP real entre contêineres, o que torna o comportamento observável de ponta a ponta.

Os endpoints principais dos nós incluem:

\begin{itemize}
    \item \code{GET /health} para verificação básica;
    \item \code{GET /state} para inspecionar o estado local;
    \item \code{GET /events} para recuperar eventos registrados;
    \item \code{POST /messages} para receber mensagens entre nós;
    \item \code{POST /commands/local-event} para gerar evento local;
    \item \code{POST /commands/send-message} para enviar mensagem a outro nó;
    \item \code{POST /commands/request-critical-section} para solicitar a seção crítica;
    \item \code{POST /commands/start-election} para disparar uma eleição manual.
\end{itemize}